\setcounter{ex}{0}

% --- CÂU 1 ---
\begin{ex}
	Gọi $M, m$ lần lượt là giá trị lớn nhất, giá trị nhỏ nhất của hàm số $y = 3\sin x + 4\cos x + 1$. Tính $M + m$.
	\shortans{2}
	\loigiai{
		Áp dụng bất đẳng thức Bunyakovsky (hoặc biến đổi về dạng $\sqrt{a^2 + b^2}\sin(x + \varphi)$), ta có:
		$$-\sqrt{3^2 + 4^2} \le 3\sin x + 4\cos x \le \sqrt{3^2 + 4^2} \iff -5 \le 3\sin x + 4\cos x \le 5 \implies -4 \le y \le 6.$$
		Do đó $M = 6, m = -4 \implies M + m = 6 + (-4) = 2$.
	}
\end{ex}

% --- CÂU 2 ---
\begin{ex}
	Cho $\tan\alpha = 3$. Tính $\cos 2\alpha$ (viết kết quả dạng số thập phân).
	\shortans{-0,8}
	\loigiai{
		Ta có:
		$$\cos 2\alpha = \dfrac{1 - \tan^2\alpha}{1 + \tan^2\alpha} = \dfrac{1 - 3^2}{1 + 3^2} = -\dfrac{8}{10} = -0{,}8.$$
	}
\end{ex}

% --- CÂU 3 ---
\begin{ex}
	Bạn Lan tiết kiệm: tháng đầu để dành $2$ triệu đồng, mỗi tháng sau để dành nhiều hơn tháng liền trước $0{,}3$ triệu đồng. Sau $12$ tháng bạn Lan tiết kiệm được bao nhiêu triệu đồng?
	\shortans{43,8}
	\loigiai{
		Số tiền tiết kiệm mỗi tháng lập thành một cấp số cộng có số hạng đầu $u_1 = 2$ (triệu đồng), công sai $d = 0{,}3$ (triệu đồng) và $n = 12$.\\
		Tổng số tiền bạn Lan tiết kiệm được sau $12$ tháng là:
		$$S_{12} = 12u_1 + \dfrac{12 \cdot 11}{2}d = 12 \cdot 2 + 66 \cdot 0{,}3 = 24 + 19{,}8 = 43{,}8\text{ (triệu đồng)}.$$
	}
\end{ex}

% --- CÂU 4 ---
\begin{ex}
	Nhiệt độ ($^\circ\text{C}$) tại một địa điểm lúc $t$ giờ ($0 \le t \le 24$) là $T(t) = 28 + 4\cos\left(\dfrac{\pi}{12}(t - 14)\right)$. Lần đầu tiên trong ngày nhiệt độ đạt $30^\circ\text{C}$ là lúc mấy giờ?
	\shortans{10}
	\loigiai{
		Khi nhiệt độ đạt $30^\circ\text{C}$, ta có:
		$$28 + 4\cos\left(\dfrac{\pi}{12}(t - 14)\right) = 30 \iff \cos\left(\dfrac{\pi}{12}(t - 14)\right) = \dfrac{1}{2}$$
		$$\iff \dfrac{\pi}{12}(t - 14) = \pm\dfrac{\pi}{3} + k2\pi \iff t = 14 \pm 4 + 24k\ (k \in \mathbb{Z}).$$
		Vì $0 \le t \le 24$ nên ta tìm được hai thời điểm là $t = 10$ và $t = 18$ (ứng với $k = 0$).\\
		Vậy lần đầu tiên trong ngày nhiệt độ đạt $30^\circ\text{C}$ là lúc $10$ giờ.
	}
\end{ex}

% --- CÂU 5 ---
\begin{ex}
	Cho hình chóp $S.ABCD$ có đáy $ABCD$ là hình thang ($AB \parallel CD,\ AB = 2CD$), $O = AC \cap BD$, $M$ là trung điểm của $SC$, $I = AM \cap SO$. Biết $\dfrac{SI}{SO} = \dfrac{a}{b}$ với $\dfrac{a}{b}$ là phân số tối giản. Tính $a + b$.
	\shortans{8}
	\loigiai{
		Trong hình thang $ABCD$ ($AB \parallel CD,\ AB = 2CD$), theo định lí Thales ta có:
		$$\dfrac{OA}{OC} = \dfrac{AB}{CD} = 2 \implies \dfrac{OC}{AC} = \dfrac{1}{3}.$$
		Trong tam giác $SAC$, kẻ $OK \parallel AM\ (K \in SC)$. Theo định lí Thales:
		\begin{itemize}
			\item Trong tam giác $AMC$: $\dfrac{CK}{CM} = \dfrac{CO}{CA} = \dfrac{1}{3} \implies MK = \dfrac{2}{3}CM = \dfrac{2}{3}SM$ (do $M$ là trung điểm $SC$).
			\item Trong tam giác $SOK$: $\dfrac{SI}{SO} = \dfrac{SM}{SK} = \dfrac{SM}{SM + MK} = \dfrac{1}{1 + \frac{2}{3}} = \dfrac{3}{5}$.
		\end{itemize}
		Vì $\dfrac{3}{5}$ là phân số tối giản nên $a = 3, b = 5 \implies a + b = 3 + 5 = 8$.
	}
\end{ex}

% --- CÂU 6 ---
\begin{ex}
	Mực nước thủy triều (mét) tại thời điểm $t$ (giờ) được cho bởi $h(t) = 6 + 2{,}5\sin\left(\dfrac{\pi t}{12} - \dfrac{\pi}{6}\right)$. Tính mực nước thấp nhất (đơn vị: mét).
	\shortans{3,5}
	\loigiai{
		Vì $\sin\left(\dfrac{\pi t}{12} - \dfrac{\pi}{6}\right) \ge -1, \forall t \ge 0$ nên:
		$$h(t) = 6 + 2{,}5\sin\left(\dfrac{\pi t}{12} - \dfrac{\pi}{6}\right) \ge 6 - 2{,}5 = 3{,}5\text{ (mét)}.$$
		Vậy mực nước thấp nhất là $3{,}5$ mét.
	}
\end{ex}
